\chapter{CONCLUSIONES Y TRABAJO FUTURO}
\label{cap:Conclusiones}
\markboth{\thechapter. Conclusiones}{\thechapter. Conclusiones}

Los resultados obtenidos en el área del control de la plataforma Stewart demuestran distintas ventajas que caracterizan a los diferentes controladores propuestos. Por un lado los controladores PID descentralizados diseñados presentan una implementación más sencilla en el diseño de control, pero carecen de rendimiento en comparación a la versión PID MIMO full la cual considera todas las relaciones cruzadas en el control de los brazos de la plataforma. Esta mejora significativa del rendimiento con el controlador PID MIMO se ve reflejado al reducir considerablemente los errores de seguimiento y al acotar las fuerzas de actuación en los brazos de la plataforma. Esto se evidencia en las distintas gráficas obtenidas y en las tablas comparativas de los distintos índices de desempeño analizados. La importancia de la diferencia en los errores y en la fuerza aplicada a los actuadores, junto a la complejidad de la implementación del control, constituye un criterio de diseño para elegir entre un controlador PID MIMO full o un controlador PID descentralizado.

Por otro lado, el seguimiento de señales sinusoidales en la plataforma Stewart agrega otro criterio en el diseño de control, al necesitar específicamente de reguladores con feedforward de la referencia para su realización. El incluir el concepto de Output Regulation en el diseño de control logra garantizar mejoras evidentes en el seguimiento sinusoidal, al bajar significativamente los errores de seguimiento. Este concepto en combinación con una acción integral, permite seguir de forma correcta a cualquier señal sinusoidal desplazada (sinusoidal sumada a una señal tipo escalón) al eliminar toda componente constante en los errores de seguimiento. Con las gráficas obtenidas y las tablas comparativas de los índices de desempeño analizados, es posible concluir que un control LQI con Output Regulation es la mejor opción para referencias sinusoidales. Este control garantiza un error en estado estacionario muy pequeño, aunque no cero, debido principalmente a las no-linealidades del sistema y las limitaciones físicas de trabajar directamente con la posición de los brazos sin el bloque de cinemática inversa (como se discutió al final del capitulo \ref{cap:OutputRegulation}). Aun así, logra ser muy eficiente en comparación con los controladores MIMO PID Full y PID descentralizado, que no están diseñados para el seguimiento de señales sinusoidal. \\

Como trabajo futuro, se buscará aplicar el bloque de cinemática inversa que transforma el movimiento de la plataforma superior al movimiento de los brazos, junto a un bloque detector de amplitud, frecuencia y fase personalizado para cada brazo actuador, que junto al diseño de control de seguimiento a señales sinusoidales formarán el esquema de control completo de la plataforma Stewart. Además, los diseños de control desarrollados en la presente Tesis se implementarán en una plataforma Stewart real, la cual sera construida y probada con un equipo de trabajo. 
